\begin{ejercicio}
    % ID: 1
    Mostrar que $f_*$ está bien definida y satisface que
    \begin{itemize}[label=(\arabic*)]
        \item [(1)] ${\id_X}_*=\id_{\pi_0(X)}.$
        \item [(2)] Para cualquier $f:X\to Y$ y $g:Y\to Z$ funciones continuas, entonces $(g\circ f)_*=g_*\circ f_*$
    \end{itemize}
%\begin{dem}
%    Supongamos $x\simeq y$, entonces existe una trayectoría continua $\sigma:I\to X$ tal que $\sigma(0)=x$ y $\sigma(1)=y$. Entonces $f\circ \sigma(0)=f(x)$ y $f\circ \sigma(1)=f(y)$
%\end{dem}
\end{ejercicio}

\begin{demostracion}
Supongamos que $x\simeq y$, entonces existe una trayectoría continua $\sigma:I\to X$ tal que $\sigma(0)=x$ y $\sigma(1)=y$. Entonces $f\circ \sigma(0)=f(x)$ y $f\circ \sigma(1)=f(y)$. Se sigue que $f\circ \sigma: I\to Y$ es testìgo de que $f(x)\simeq f(y)$, por lo tanto, 
$$f_*\langle x\rangle=\langle f(x)\rangle=\langle f(y)\rangle=f_*\langle y\rangle,$$
i.e., $f_*$ está bien definida. 

Hagamos $x\in X$,
\begin{enumerate}[label=(\arabic*)]
    \item Entonces 
    \begin{align*}
        {\id_X}_*\langle x\rangle &=\langle\id_X(x)\rangle=\langle x\rangle,
    \end{align*}
    por la tanto, ${\id_X}_*=\id_{\pi_0(X)}.$
    \item Por otro lado,
    \begin{align*}
        (g\circ f)_*\langle x\rangle
        &=\langle g\circ f(x)\rangle\\
        &=\langle g(f(x))\rangle\\
        &=g_*\langle f(x)\langle\\
        &=g_*(f_*\langle x\rangle),
    \end{align*}
    por lo tanto, $(g\circ f)_*=g_*\circ f_*$.
\end{enumerate}
\end{demostracion}
